\chapter{Sprint 0: Requirements Specification}

\section*{Introduction}
\addcontentsline{toc}{section}{Introduction}

Sprint 0 is a discovery and preparation sprint. This chapter identifies the system actors, lists
the functional and non-functional requirements, organizes them into a product backlog, describes
the work tools adopted by the team, and defines the general architecture together with static and
dynamic design views that guide the implementation sprints.

%% Clickable back-reference to a webography item: \webref{web:fastapi}
\newcommand{\webref}[1]{\hyperref[#1]{[ref.~\ref*{#1}]}}

\newcommand{\toollogo}[3][2.0cm]{%
  \begin{figure}[H]
  \centering
  \includegraphics[height=#1,keepaspectratio]{#2}\\[0.3em]
  {\small\textbf{#3}}
  \end{figure}
}

\newcommand{\toollogotriple}[6]{%
  \begin{figure}[H]
  \centering
  \begin{minipage}[t]{0.28\textwidth}
    \centering
    \includegraphics[height=1.45cm,keepaspectratio]{#1}\\[0.25em]
    {\small\textbf{#2}}
  \end{minipage}
  \hfill
  \begin{minipage}[t]{0.28\textwidth}
    \centering
    \includegraphics[height=1.45cm,keepaspectratio]{#3}\\[0.25em]
    {\small\textbf{#4}}
  \end{minipage}
  \hfill
  \begin{minipage}[t]{0.28\textwidth}
    \centering
    \includegraphics[height=1.45cm,keepaspectratio]{#5}\\[0.25em]
    {\small\textbf{#6}}
  \end{minipage}
  \end{figure}
}

\newcommand{\toollogodouble}[4]{%
  \begin{figure}[H]
  \centering
  \begin{minipage}[t]{0.38\textwidth}
    \centering
    \includegraphics[height=1.55cm,keepaspectratio]{#1}\\[0.25em]
    {\small\textbf{#2}}
  \end{minipage}
  \hfill
  \begin{minipage}[t]{0.38\textwidth}
    \centering
    \includegraphics[height=1.55cm,keepaspectratio]{#3}\\[0.25em]
    {\small\textbf{#4}}
  \end{minipage}
  \end{figure}
}

%% ─── 2.1 ─────────────────────────────────────────────────────────────────────
\section{System Actors}

The BVIRAL Platform has four actor types:

\begin{table}[H]
\centering
\caption{System actors and their roles}
\begin{tabularx}{\textwidth}{lX}
\toprule
\textbf{Actor} & \textbf{Description} \\
\midrule
\textbf{Anonymous Visitor} & Any unauthenticated user who can browse the video portal, search
  for clips, and interact with the chatbot assistant. \\[0.4em]
\textbf{Auth. Creator / Partner} & A registered user who has completed onboarding and holds an
  active subscription. Can access the dashboard, manage playlists, submit videos, and use
  the billing portal. \\[0.4em]
\textbf{Administrator} & Internal BVIRAL staff with elevated access to review custom quote
  requests, manage pricing, and view admin reports. \\[0.4em]
\textbf{External System} & Stripe (payment webhooks), Clerk (authentication tokens), DAM
  provider APIs (media access provisioning), and the SMS Scraper API (social media follower data). \\
\bottomrule
\end{tabularx}
\end{table}

%% ─── 2.2 ─────────────────────────────────────────────────────────────────────
\section{Functional Requirements}

\subsection{Video Search Portal}

\begin{enumerate}[label=\textbf{FR-VP-\arabic*}]
  \item Browse the video catalogue with pagination and category filtering.
  \item Search videos by keyword or semantic query.
  \item View detailed video information: title, duration, category, resolution, tags.
  \item Play video previews directly in the browser.
  \item Access an embedded AI chatbot for natural language queries.
  \item Transcribe spoken audio queries via the microphone.
  \item View the latest and trending clips on the homepage.
  \item View a personalized \textbf{Recommended for You} section on the homepage, driven by
    the user's past searches and click events.
  \item Save videos to favorites and remove them; view the full favorites list.
\end{enumerate}

\subsection{Authentic Rights Portal (Onboarding and Dashboard)}

\begin{enumerate}[label=\textbf{FR-AR-\arabic*}]
  \item Submit social media channel URLs (Instagram, TikTok, YouTube, Facebook, Twitter/X,
    Snapchat) to begin onboarding.
  \item Receive automatic follower count analysis; channels over 2M followers trigger a
    custom pricing flow.
  \item Select a monthly or yearly subscription plan and complete Stripe embedded checkout.
  \item Sign the digital service agreement before checkout; receive a signed PDF by email.
  \item Complete account setup (name, account type, password) after successful payment.
  \item Access a personal dashboard with subscription status, renewal alerts, and billing portal.
  \item Manage playlists: create, rename, add or remove videos.
  \item Browse and search the licensed video \textbf{library}: filter by keyword and category,
    view video details, and download licensed clips.
  \item Manage \textbf{favorites}: add or remove videos and view the full favorites list.
  \item Submit new video clips for licensing review.
  \item Interact with the embedded AI support chatbot.
  \item \textit{Administrator}: review pending custom quote requests and set pricing.
  \item \textit{Administrator}: access whitelist reports and subscription announcements.
  \item \textit{Administrator}: browse and inspect all user chat conversations.
  \item \textit{Administrator}: test the AI chatbot directly from an admin console.
  \item \textit{Administrator}: manage the FAQ knowledge base — add, edit, search, and delete
    FAQ entries ingested into the RAG knowledge base.
  \item \textit{Administrator}: review submitted video clips and update their status.
\end{enumerate}

\subsection{Agentic Helper}

\begin{enumerate}[label=\textbf{FR-AH-\arabic*}]
  \item Accept chat messages and route them by intent (RAG, direct LLM, or video-search mode).
  \item Retrieve relevant knowledge chunks via Qdrant cosine similarity for support queries.
  \item Extract structured search filters (category, tags, duration, resolution, orientation)
    from free-text queries and execute a \textbf{deterministic advanced search} against the
    Videos Search API as the primary video-search path.
  \item Fall back to the \textbf{ReAct video-search agent} when the deterministic
    search service is unavailable.
  \item Return video results with a search URL for direct portal navigation.
  \item Ingest knowledge documents (JSON or Excel) into the RAG knowledge base.
  \item Add, update, search, and delete individual FAQ entries in the RAG knowledge base via
    dedicated admin endpoints.
  \item Transcribe audio input using OpenAI Whisper.
  \item Maintain per-user conversation history with a configurable memory window.
\end{enumerate}

%% ─── 2.3 ─────────────────────────────────────────────────────────────────────
\section{Non-Functional Requirements}

\begin{table}[H]
\centering
\caption{Non-functional requirements}
\begin{tabularx}{\textwidth}{lX}
\toprule
\textbf{Category} & \textbf{Requirement} \\
\midrule
\textbf{Performance} & Video search returns results within 500\,ms at the 95th percentile.
  Chat responses begin streaming within 2 seconds. \\[0.3em]
\textbf{Scalability} & Each microservice scales independently via container replicas.
  OpenSearch supports horizontal shard scaling. Celery workers add capacity on demand. \\[0.3em]
\textbf{Availability} & Target uptime of 99.5\% in production. Services expose \texttt{/health}
  and \texttt{/health} endpoint for load-balancer probes. \\[0.3em]
\textbf{Security} & All endpoints require authentication (Clerk JWT Bearer or API key). Stripe
  webhook payloads are verified with HMAC. Secrets are stored in per-service \texttt{.env}
  files, excluded from version control via \texttt{.gitignore}; nothing sensitive is committed
  to the repository. \\[0.3em]
\textbf{Maintainability} & Domain-driven module structure (\texttt{router → services → models}).
  Alembic manages all schema changes. Prompts are versioned text files. \\[0.3em]
\textbf{Observability} & Sentry error tracking in production for all Python services.
  All Celery tasks include retry/backoff metadata. \\[0.3em]
\textbf{Portability} & The full stack runs with a single \texttt{docker compose up} command, with no manual environment setup required. \\
\bottomrule
\end{tabularx}
\end{table}

%% ─── 2.4 ─────────────────────────────────────────────────────────────────────
\section{Use-Case Diagrams}

\subsection{Global Use-Case Diagram}

Figure~\ref{fig:uc_global} shows all system use cases grouped by portal, with the human and
system actors that interact with each.

\begin{figure}[H]
\centering
\includegraphics[width=\textwidth]{\detokenize{diagrams/global use.drawio.png}}
\caption{Global use-case diagram of the BVIRAL Platform}
\label{fig:uc_global}
\end{figure}

%% ─── 2.5 ─────────────────────────────────────────────────────────────────────
\section{Product Backlog}

The product backlog groups the functional requirements into user stories organized per epic and
prioritized with the industrial supervisor. Priority levels follow the MoSCoW convention
(\textbf{M}ust have, \textbf{S}hould have, \textbf{C}ould have).

\small
\begin{longtable}{@{}llp{0.52\textwidth}c@{}}
\caption{Product backlog (user stories per epic)}
\label{tab:product_backlog} \\
\toprule
\textbf{ID} & \textbf{Epic} & \textbf{User story} & \textbf{Prio.} \\
\midrule
\endfirsthead
\multicolumn{4}{c}{\tablename~\thetable{} — \textit{continued}} \\
\toprule
\textbf{ID} & \textbf{Epic} & \textbf{User story} & \textbf{Prio.} \\
\midrule
\endhead
\midrule
\multicolumn{4}{r}{\textit{Continued on next page}} \\
\endfoot
\bottomrule
\endlastfoot
\multicolumn{4}{l}{\textit{Sprint 1 — Authentication, Onboarding and Payments}} \\
\midrule
US-01 & Authentication  & As a visitor, I want to sign in through Clerk so that I can access
  the rights portal. & M \\
US-02 & Onboarding      & As a creator, I want to submit my social channels so the platform can
  analyze my reach. & M \\
US-03 & Onboarding      & As a creator with $>$2M followers, I want a custom-quote flow so that
  pricing can be negotiated. & S \\
US-04 & Payments        & As a user, I want to complete Stripe embedded checkout and sign the
  service agreement before payment. & M \\
US-05 & Notifications   & As a user, I want to receive a confirmation email with a signed PDF
  once payment is captured. & M \\
US-06 & Dashboard       & As a subscriber, I want a dashboard showing subscription status,
  renewal date, and a billing-portal link. & M \\
\midrule
\multicolumn{4}{l}{\textit{Sprint 2 — Video Indexing and Keyword Search}} \\
\midrule
US-07 & Video search    & As a visitor, I want keyword-based video search with filters
  (category, duration, resolution). & M \\
US-08 & Video search    & As a visitor, I want to browse the catalogue by category with
  pagination. & M \\
\midrule
\multicolumn{4}{l}{\textit{Sprint 3 — Semantic and Parallel Search with Ranking}} \\
\midrule
US-09 & Video search    & As a visitor, I want semantic search so natural-language queries work
  even without matching keywords. & M \\
US-10 & Video search    & As a visitor, I want a hybrid ranking that combines lexical and
  semantic signals for better relevance. & S \\
\midrule
\multicolumn{4}{l}{\textit{Sprint 4 — Agentic Helper: RAG and Classifier Pipeline}} \\
\midrule
US-11 & Agentic AI      & As a visitor, I want a chatbot that understands my query and returns
  either support answers or video results. & M \\
US-12 & Agentic AI      & As a support user, I want answers grounded in the BVIRAL knowledge
  base via RAG. & M \\
\midrule
\multicolumn{4}{l}{\textit{Sprint 5 — Search Orchestrator, Voice Transcription and Portals}} \\
\midrule
US-13 & Agentic AI      & As a user, I want to dictate my query with a microphone and have it
  transcribed into the chat input. & S \\
US-14 & Playlists       & As a subscriber, I want to create and manage playlists of licensed
  videos. & S \\
US-15 & Library         & As a subscriber, I want to browse my licensed video library, search
  by keyword or category, and download clips. & M \\
US-16 & Video submission & As a subscriber, I want to submit new video clips for licensing review
  and track the status of my submissions. & S \\
US-17 & Admin           & As an admin, I want to review custom-quote requests, set pricing, and
  view the status history of all quotes. & S \\
US-18 & Admin           & As an admin, I want to review submitted video clips and update their
  approval status. & S \\
US-19 & Admin           & As an admin, I want to inspect user chat conversations, test the
  chatbot from an admin console, and manage FAQ knowledge-base entries. & S \\
\midrule
\multicolumn{4}{l}{\textit{Sprint 6 — Personalized Recommendations}} \\
\midrule
US-20 & Personalization & As a subscriber, I want the platform to learn from my searches and
  clicks to surface relevant videos in my feed. & S \\
US-21 & Favorites       & As a subscriber, I want to save videos to a favorites list and remove
  them at any time. & S \\
\end{longtable}
\normalsize

%% ─── 2.6 ─────────────────────────────────────────────────────────────────────
\section{Work Tools}

\subsection{Frameworks and Libraries}

\textbf{FastAPI}

FastAPI is used to build the AR-API and Agentic Helper backend services with asynchronous endpoints,
automatic OpenAPI documentation, and clear request/response schema validation~\webref{web:fastapi}.

\toollogo{figures/logos/fastapilogo.png}{FastAPI}

\textbf{Flask + Flask-RESTX}

Flask with Flask-RESTX is used in the Videos Search API to expose REST endpoints and organize service
resources around OpenSearch video indexing and query operations~\webref{web:flask}.

\toollogo{figures/logos/flasklogo.png}{Flask}

\textbf{SQLAlchemy 2.0 async + asyncpg}

SQLAlchemy asynchronous ORM with asyncpg provides non-blocking PostgreSQL data access in FastAPI
services, enabling efficient database interactions under concurrent API workloads~\webref{web:sqlalchemy}.

\toollogo{figures/logos/SQLAlchemy-logo.png}{SQLAlchemy}

\textbf{Alembic}

Alembic is used for database schema migration management, allowing controlled evolution of PostgreSQL
tables, constraints, and relationships across project iterations~\webref{web:alembic}.

\toollogo{figures/logos/alembic logo.png}{Alembic}

\textbf{Celery + Redis}

Celery with Redis is used to process asynchronous and scheduled background tasks, such as delayed
operations and periodic jobs that should not block API request handling~\webref{web:celery}\webref{web:redis}.

\toollogodouble
  {figures/logos/celery logo.png}{Celery}
  {figures/logos/Logo-redis.png}{Redis}

\textbf{LangChain}

LangChain is used to implement the chatbot orchestration logic, including prompt pipelines, tool
calling behavior, and multi-step reasoning flows for video search interactions~\webref{web:langchain}.

\toollogo{figures/logos/LangChain_Logo.png}{LangChain}

\textbf{Qdrant}

Qdrant is the vector database used to store embeddings and perform semantic similarity retrieval,
supporting context grounding for the AI assistant knowledge workflow~\webref{web:qdrant}.

\toollogo{figures/logos/qdrant-logo.png}{Qdrant}

\textbf{OpenSearch Python}

The OpenSearch Python client is used by the Videos Search API to query and manage indexed video
documents with lexical search and vector-aware retrieval capabilities~\webref{web:opensearch}.

\toollogo{figures/logos/opensearch logo.png}{OpenSearch}

\textbf{Next.js 15 / 16}

Next.js powers both web portals and provides server-side rendering, App Router features, and API
routes used as backend-for-frontend layers for integration with internal services~\webref{web:nextjs}.

\toollogo{figures/logos/Nextjs-logo.png}{Next.js}

\textbf{Radix UI / Material UI / Tailwind CSS}

These UI libraries are used to accelerate frontend implementation with reusable components, theming,
and utility-first styling for consistent and responsive interfaces~\webref{web:radixui}\webref{web:mui}\webref{web:tailwind}.

\toollogotriple
  {figures/logos/radix ui logo.png}{Radix UI}
  {figures/logos/mui logo.png}{Material UI}
  {figures/logos/tailwind logo.png}{Tailwind CSS}

\textbf{Clerk}

Clerk is used for authentication and user identity management, including secure sign-in flows and
session handling across the platform portals and protected backend endpoints~\webref{web:clerk}.

\toollogo{figures/logos/clerk logo.png}{Clerk}

\textbf{Stripe JS + Python}

Stripe SDKs are used for subscription billing workflows, including checkout session handling, payment
integration, and secure webhook verification in backend services~\webref{web:stripe}.

\toollogo{figures/logos/stripe logo.png}{Stripe}

\subsection{Programming Languages}

\begin{itemize}
  \item \textbf{Python 3.12} for AR-API and Agentic Helper.
  \item \textbf{Python 3.8} for Videos Search API (compatibility with the OpenSearch Python client).
  \item \textbf{TypeScript} (Node.js 20 LTS) for both Next.js portals.
  \item \textbf{SQL} for migrations (PostgreSQL 15/16) and OpenSearch query DSL (JSON).
\end{itemize}

\subsection{Tools Used}

\textbf{Cursor IDE}

Cursor IDE is the main development environment used in this project. It provides an AI-assisted
workflow for writing, refactoring, and reviewing code across the platform services~\webref{web:cursor}.

\toollogo{figures/logos/cursor ide logo.final.png}{Cursor IDE}

\textbf{Docker}

Docker is used to package services in isolated containers so each component runs with consistent
runtime dependencies across development environments~\webref{web:docker}.

\toollogo{figures/logos/Docker_Logo.png}{Docker}

\textbf{Docker Compose}

Docker Compose orchestrates multi-service local environments by defining and starting all required
containers (APIs, databases, cache, and supporting services) from a single configuration~\webref{web:dockercompose}.

\toollogo{figures/logos/docker_compose_logo.jpeg}{Docker Compose}

\textbf{Git}

Git is used for distributed version control, enabling feature branching, history tracking, and safe
collaboration during iterative development.

\toollogo{figures/logos/Git-logo.png}{Git}

\textbf{Postman}

Postman is used to test and validate REST API endpoints by building requests, inspecting responses,
and organizing reusable API collections~\webref{web:postman}.

\toollogo{figures/logos/Postman_logo.png}{Postman}

\textbf{Qdrant Cloud Console}

Qdrant Cloud Console is used to inspect vector collections, monitor indexed embeddings, and verify
semantic retrieval data used by the chatbot pipeline~\webref{web:qdrant_cloud_console}.

\toollogo{figures/logos/qdrant cloud console logo.png}{Qdrant Cloud Console}

\textbf{OpenSearch Dashboards}

OpenSearch Dashboards provides a web interface to explore indexed video data, validate mappings, and
analyze search behavior in the OpenSearch engine~\webref{web:opensearch_dashboards}.

\toollogo{figures/logos/opensearch logo.png}{OpenSearch Dashboards}

%% ─── 2.7 ─────────────────────────────────────────────────────────────────────
\section{General Architecture}

The BVIRAL Platform uses a \textbf{microservices architecture}. Each service is independently
deployable, has a defined API contract, and owns its own data store. Services communicate over
HTTP REST, and there is no shared database.

\begin{figure}[H]
\centering
\resizebox{\textwidth}{!}{%
\begin{tikzpicture}[x=1cm, y=1cm,
  %% ── Node styles ──────────────────────────────────────────────────────────
  %% fe  : min-width 5.5 cm  → half = 2.75 cm
  %% api : min-width 4.5 cm  → half = 2.25 cm
  %% db  : min-width 3.0 cm  → half = 1.50 cm
  %% ext : min-width 3.2 cm  → half = 1.60 cm
  fe/.style={draw=blue!65!black, fill=blue!10, rounded corners=6pt,
             minimum width=5.5cm, minimum height=1.15cm,
             font=\small\bfseries, align=center, text=blue!15!black},
  api/.style={draw=teal!65!black, fill=teal!10, rounded corners=6pt,
              minimum width=4.5cm, minimum height=1.15cm,
              font=\small\bfseries, align=center, text=teal!15!black},
  db/.style={draw=violet!60, fill=violet!8, rounded corners=6pt,
             minimum width=3.0cm, minimum height=0.95cm,
             font=\small\bfseries, align=center},
  ext/.style={draw=gray!55, fill=gray!8, rounded corners=6pt,
              minimum width=3.2cm, minimum height=0.95cm,
              font=\small\bfseries, align=center},
  %% ── Arrow styles ─────────────────────────────────────────────────────────
  arr/.style={->, >=Stealth, thick, gray!65},
  extarr/.style={->, >=Stealth, semithick, dashed, gray!50}
]

%% ── ROW 1  Client Layer  (y = 8.5) ──────────────────────────────────────────
%% VSP  x=[-6.75, -1.25]   ARP  x=[1.25, 6.75]   gap = 2.5 cm
\node[fe] (vp)  at (-4.0, 8.5) {Video Search Portal\\{\scriptsize Next.js 15 · Tailwind}};
\node[fe] (arp) at ( 4.0, 8.5) {Authentic Rights Portal\\{\scriptsize Next.js 16 · MUI}};

%% ── ROW 2  Backend Services  (y = 5.5) ──────────────────────────────────────
%% Ahai x=[-9.25,-4.75]  VSAPI x=[-2.25,2.25]  ARAPI x=[4.75,9.25]  gap=2.5cm
\node[api] (ahai)  at (-7.0, 5.5) {Agentic Helper\\{\scriptsize FastAPI · :8003}};
\node[api] (vsapi) at ( 0.0, 5.5) {Videos Search API\\{\scriptsize Flask · :5000}};
\node[api] (arapi) at ( 7.0, 5.5) {AR-API\\{\scriptsize FastAPI · :8001}};

%% ── ROW 3  Data Layer  (y = 2.5) ────────────────────────────────────────────
%% Qdrant x=[-8.5,-5.5]  OS x=[-2.0,1.0]  PG x=[2.0,5.0]  RD x=[6.5,9.5]
%% gaps:  Qdrant–OS = 3.5 cm   OS–PG = 1.0 cm   PG–RD = 1.5 cm
\node[db] (qdrant) at (-7.0, 2.5) {Qdrant Cloud\\{\scriptsize :6333}};
\node[db] (os)     at (-0.5, 2.5) {OpenSearch\\{\scriptsize :9200}};
\node[db] (pg)     at ( 3.5, 2.5) {PostgreSQL\\{\scriptsize :5433}};
\node[db] (rd)     at ( 8.0, 2.5) {Redis\\{\scriptsize :6379}};

%% ── External Left ────────────────────────────────────────────────────────────
\node[ext] (oai) at (-11.5, 5.5) {OpenAI API};

%% ── External Right ───────────────────────────────────────────────────────────
\node[ext] (stripe) at (12.0, 6.5) {Stripe};
\node[ext] (clerk)  at (12.0, 5.5) {Clerk};
\node[ext] (dam)    at (12.0, 4.5) {DAM Provider};

%% ── Layer background boxes (manually bounded — no fit/positioning issues) ───
\begin{scope}[on background layer]
  %% Client layer  — covers fe nodes ±0.6 cm padding
  \fill[blue!3, rounded corners=8pt]
    (-7.2, 7.7) rectangle (7.2, 9.3);
  \draw[blue!30, dashed, rounded corners=8pt, line width=0.8pt]
    (-7.2, 7.7) rectangle (7.2, 9.3);
  %% Backend services — covers api nodes ±0.6 cm padding
  \fill[teal!3, rounded corners=8pt]
    (-9.7, 4.7) rectangle (9.7, 6.3);
  \draw[teal!30, dashed, rounded corners=8pt, line width=0.8pt]
    (-9.7, 4.7) rectangle (9.7, 6.3);
  %% Data layer — covers db nodes ±0.6 cm padding
  \fill[violet!3, rounded corners=8pt]
    (-9.0, 1.7) rectangle (9.9, 3.3);
  \draw[violet!25, dashed, rounded corners=8pt, line width=0.8pt]
    (-9.0, 1.7) rectangle (9.9, 3.3);
\end{scope}

%% ── Layer Labels (anchored at top-left of each box) ─────────────────────────
\node[font=\fontsize{7.5}{9}\selectfont\bfseries, text=blue!45!gray,
      anchor=south west] at (-7.0, 9.3) {\strut CLIENT LAYER};
\node[font=\fontsize{7.5}{9}\selectfont\bfseries, text=teal!45!gray,
      anchor=south west] at (-9.5, 6.3) {\strut BACKEND SERVICES};
\node[font=\fontsize{7.5}{9}\selectfont\bfseries, text=violet!50!gray,
      anchor=south west] at (-8.8, 3.3) {\strut DATA LAYER};

%% ── Arrows: Frontends → Backend APIs ─────────────────────────────────────────
\draw[arr] (vp)  -- (ahai);
\draw[arr] (vp)  -- (vsapi);
\draw[arr] (arp) -- (arapi);
%% ARP also uses the Agentic Helper (AI chatbot) — routed above backend layer
\draw[arr] (arp.west) to[out=200, in=30] (ahai.north);

%% ── Arrows: Inter-service (Agentic Helper → Videos Search API) ───────────────
\draw[arr] (ahai.east) -- (vsapi.west);

%% ── Arrows: Backend APIs → Data Stores ───────────────────────────────────────
\draw[arr] (ahai.south)       -- (qdrant.north);
\draw[arr] (vsapi.south)      -- (os.north);
\draw[arr] (arapi.south west) -- (pg.north);
\draw[arr] (arapi.south east) -- (rd.north);

%% ── Arrows: External Services (dashed) ───────────────────────────────────────
\draw[extarr] (ahai.west)  -- (oai.east);
\draw[extarr] (arapi.east) -- (stripe.west);
\draw[extarr] (arapi.east) -- (clerk.west);
\draw[extarr] (arapi.east) -- (dam.west);

\end{tikzpicture}%
}
\caption{BVIRAL Platform — global microservices architecture}
\label{fig:arch_global}
\end{figure}

%% ─── 2.8 ─────────────────────────────────────────────────────────────────────
\section{Static Design: Class Diagram}

\subsection{Global Class Diagram}

\begin{figure}[H]
\centering
\includegraphics[width=\textwidth]{\detokenize{diagrams/global class dag.drawio.png}}
\caption{Global class diagram of the BVIRAL Platform}
\label{fig:er_arapi}
\end{figure}

\section*{Conclusion}
\addcontentsline{toc}{section}{Conclusion}

Sprint 0 produced the actors, functional and non-functional requirements, a prioritized product
backlog, the selected work tools, and the general architecture together with its static view.
These artifacts form the shared foundation for all implementation sprints. The
next chapter starts Sprint 1 with authentication, onboarding, and Stripe payments.
